\documentclass{article}
\usepackage[T1]{fontenc}
\usepackage{lmodern}
\usepackage{graphicx}
\usepackage{amsmath,amssymb}
\usepackage[a4paper,margin=2cm]{geometry}
\usepackage[absolute,overlay]{textpos}
\setlength{\TPHorizModule}{1mm}
\setlength{\TPVertModule}{1mm}
\usepackage[indonesian]{babel}
\usepackage{hyperref}
\setlength{\emergencystretch}{2em}
% unit: Halaman 38

\begin{document}

% === Halaman 38 (python/native) ===

38

• Kelebihan Vigenere Cipher: huruf plainteks yang sama tidak selalu 

dienkripsi menjadi huruf cipheteks yang sama pula, bergantung huruf 

kunci yang digunakan. 

 Contoh: pada contoh di atas, huruf plainteks T dapat dienkripsi menjadi L 

atau H,  dan huruf cipherteks V dapat merepresentasikan huruf plainteks 

H, I, dan X 

• Hal di atas merupakan karakteristik dari cipher abjad-majemuk: setiap 

huruf cipherteks dapat memiliki kemungkinan banyak huruf plainteks. 

• Bandingkan dengan cipher abjad-tunggal, setiap huruf cipherteks selalu 

menggantikan huruf plainteks tertentu.

\end{document}
